\chapter{Anexo IV-Manual de Usuario}
\label{chapter:A-3-Manual de Usuario}
\thispagestyle{empty}

\newpage
\tableofcontents
\newpage

\section{Introducción}
\label{section:A-3-Introducción}


Este manual describe el uso de la \Gls{api} de servicios en la plataforma, un sistema diseñado para gestionar cursos, lecciones, transcripciones y conversaciones mediante inteligencia artificial. La \Gls{api} proporciona funcionalidades de scraping, generación de respuestas mediante RAG (Retrieval-Augmented Generation), transcripción de contenido y gestión de usuarios.



\subsection{URL Base}
\label{subsection:A-3-URL Base}

Todos los endpoints de la \Gls{api} utilizan la siguiente \acrshort{url} base:

\begin{lstlisting}[language=bash]
http://localhost:8000
\end{lstlisting}

\section{Códigos de Estado HTTP}
\label{section:A-3-Códigos de Estado HTTP}

La \Gls{api} utiliza códigos de estado\acrshort{http} estándar:

\begin{table}[h]
\centering
\begin{tabular}{|c|l|p{7cm}|}
\hline
\textbf{Código} & \textbf{Estado} & \textbf{Descripción} \\
\hline
200 & OK & Solicitud exitosa \\
\hline
201 & Created & Recurso creado exitosamente \\
\hline
400 & Bad Request & Solicitud mal formada o parámetros inválidos \\
\hline
401 & Unauthorized & Autenticación requerida o \gls{token} inválido \\
\hline
403 & Forbidden & Sin permisos para acceder al recurso \\
\hline
404 & Not Found & Recurso no encontrado \\
\hline
500 & Internal Server Error & Error interno del servidor \\
\hline
\end{tabular}
\caption{Códigos de estado \acrshort{http} comunes}
\end{table}

\section[Métodos de la api]{Métodos de la \Gls{api}}
\label{Métodos de la api}

Esta sección describe todos los métodos disponibles en la \Gls{api}, organizados por categoría funcional.

\subsection{Autenticación}
\label{subsection:A-3-Autenticación}

Para utilizar los endpoints protegidos, debe incluir el \gls{token} de acceso en el encabezado:

\begin{lstlisting}[language=bash]
curl -X GET http://localhost:8000/endpoint \
  -H "Authorization: Bearer YOUR_ACCESS_TOKEN"
\end{lstlisting}

\subsubsection{Health Check}
\label{Health Check}

\textbf{Endpoint:} \texttt{GET /health}

\textbf{Autenticación:} No requerida

\textbf{Descripción:} Verifica que el servicio está activo y funcionando correctamente.

\textbf{Códigos de Respuesta:}
\begin{itemize}
    \item \texttt{200}: Successful Response
\end{itemize}

\textbf{Ejemplo:}

\begin{lstlisting}[language=bash]
curl -X GET http://localhost:8000/health
\end{lstlisting}

\subsubsection{Registro de usuario}
\label{subsubsection:A-3-Registro de usuario}

\textbf{Endpoint:} \texttt{POST /register}

\textbf{Autenticación:} No requerida

\textbf{Descripción:} Crea una nueva cuenta de usuario con email, nombre de usuario y contraseña. Todos los campos son obligatorios y el username/email deben ser únicos.

\textbf{Códigos de Respuesta:}
\begin{itemize}
    \item \texttt{201}: Successful Response
    \item \texttt{422}: Validation Error
\end{itemize}

\textbf{Ejemplo:}

\begin{lstlisting}[language=bash]
curl -X POST http://localhost:8000/register \
  -H "Content-Type: application/json" \
  -d '{
    "email": "usuario@example.com",
    "username": "mi_usuario",
    "password": "mi_password_seguro"
  }'
\end{lstlisting}

\subsubsection{Inicio de sesión}
\label{subsubsection:Inicio de sesión}

\textbf{Endpoint:} \texttt{POST /login}

\textbf{Autenticación:} No requerida

\textbf{Descripción:} Inicia sesión con username/password para recibir un \gls{token} JWT de acceso para solicitudes autenticadas.

\textbf{Códigos de Respuesta:}
\begin{itemize}
    \item \texttt{200}: Successful Response
    \item \texttt{422}: Validation Error
\end{itemize}

\textbf{Ejemplo:}

\begin{lstlisting}[language=bash]
curl -X POST http://localhost:8000/login \
  -H "Content-Type: application/json" \
  -d '{
    "username": "mi_usuario",
    "password": "mi_password_seguro"
  }'
\end{lstlisting}

\textbf{Respuesta esperada:}

\begin{lstlisting}[language=json]
{
  "access_token": "eyJhbGciOiJIUzI1NiIsInR5cCI6IkpXVCJ9...",
  "refresh_token": "eyJhbGciOiJIUzI1NiIsInR5cCI6IkpXVCJ9...",
  "token_type": "bearer"
}
\end{lstlisting}

\subsubsection{Renovación de token}
\label{subsubsection:Renovación de token}

\textbf{Endpoint:} \texttt{POST /refresh}

\textbf{Autenticación:} No requerida (utiliza refresh \gls{token})

\textbf{Descripción:} Genera un nuevo \gls{token} de acceso a partir del refresh \gls{token}.

\textbf{Códigos de Respuesta:}
\begin{itemize}
    \item \texttt{200}: Successful Response
    \item \texttt{422}: Validation Error
\end{itemize}

\textbf{Ejemplo:}

\begin{lstlisting}[language=bash]
curl -X POST http://localhost:8000/refresh \
  -H "Content-Type: application/json" \
  -d '{
    "refresh_token": "eyJhbGciOiJIUzI1NiIsInR5cCI6IkpXVCJ9..."
  }'
\end{lstlisting}

\subsection{Servicios de Scraper}
\label{subsection:A-3-Servicios de Scraper}

\subsubsection{Sincronizar todos los cursos}
\label{subsubsection:A-3-Sincronizar todos los cursos}

\textbf{Endpoint:} \texttt{POST /v1/services/scraper/sync-all}

\textbf{Autenticación:} Requerida

\textbf{Descripción:} Realiza scraping de todos los cursos y lecciones desde \Gls{OpenFing} y los sincroniza con la base de datos.

\textbf{Códigos de Respuesta:}
\begin{itemize}
    \item \texttt{200}: Successful Response
\end{itemize}

\textbf{Ejemplo:}

\begin{lstlisting}[language=bash]
curl -X POST http://localhost:8000/v1/services/scraper/sync-all \
  -H "Authorization: Bearer YOUR_ACCESS_TOKEN"
\end{lstlisting}

\subsubsection{Sincronizar un curso específico}
\label{subsubsection:A-3-Sincronizar un curso específico}

\textbf{Endpoint:} \texttt{POST /v1/services/scraper/sync-course/\{course\_code\}}

\textbf{Autenticación:} Requerida

\textbf{Parámetros de ruta:}
\begin{itemize}
    \item \texttt{course\_code} (string, requerido): Código del curso a sincronizar
\end{itemize}

\textbf{Descripción:} Realiza scraping de un curso específico y sus lecciones desde \Gls{OpenFing}.

\textbf{Códigos de Respuesta:}
\begin{itemize}
    \item \texttt{200}: Successful Response
    \item \texttt{422}: Validation Error
\end{itemize}

\textbf{Ejemplo:}

\begin{lstlisting}[language=bash]
curl -X POST http://localhost:8000/v1/services/scraper/sync-course/PROG1 \
  -H "Authorization: Bearer YOUR_ACCESS_TOKEN"
\end{lstlisting}

\subsection{Servicios de Generación (Chat IA)}
\label{subsection:A-3-Servicios de Generación (Chat IA)}

\subsubsection{Enviar mensaje al chat}
\label{subsubsection:A-3-Enviar mensaje al chat}

\textbf{Endpoint:} \texttt{POST /v1/services/generator/chat}

\textbf{Autenticación:} Requerida

\textbf{Descripción:} Procesa un mensaje de chat utilizando RAG (Retrieval-Augmented Generation). Soporta usuarios anónimos (efímeros), usuarios JWT y usuarios con \Gls{api} key.

\textbf{Códigos de Respuesta:}
\begin{itemize}
    \item \texttt{200}: Successful Response
    \item \texttt{422}: Validation Error
\end{itemize}

\textbf{Ejemplo:}

\begin{lstlisting}[language=bash]
DATA='{
    "message": "Explícame qué es una variable en programación",
    "conversation_id": 1,
    "course_code": "PROG1"
  }'
curl -X POST http://localhost:8000/v1/services/generator/chat \
  -H "Authorization: Bearer YOUR_ACCESS_TOKEN" \
  -H "Content-Type: application/json" \
  -d "\$DATA"
\end{lstlisting}

\subsubsection{Enviar feedback sobre un mensaje}
\label{subsubsection:A-3-Enviar feedback sobre un mensaje}

\textbf{Endpoint:} \texttt{POST /v1/services/generator/message/\{message\_id\}/feedback}

\textbf{Autenticación:} Requerida

\textbf{Parámetros de ruta:}
\begin{itemize}
    \item \texttt{message\_id} (integer, requerido): ID del mensaje
\end{itemize}

\textbf{Descripción:} Permite al usuario proporcionar feedback sobre un mensaje del sistema.

\textbf{Códigos de Respuesta:}
\begin{itemize}
    \item \texttt{200}: Successful Response
    \item \texttt{422}: Validation Error
\end{itemize}

\textbf{Ejemplo:}

\begin{lstlisting}[language=json]
curl -X POST http://localhost:8000/v1/services/generator/message/123/feedback \
  -H "Authorization: Bearer YOUR_ACCESS_TOKEN" \
  -H "Content-Type: application/json" \
  -d '{
    "rating": 5,
    "comment": "Respuesta muy {\'u}til y clara"
  }'
\end{lstlisting}

\subsection{Servicios de Transcripción}
\label{subsection:A-3-Servicios de Transcripción}


\subsubsection{Crear o actualizar transcripción}
\label{subsubsection:A-3-Crear o actualizar transcripción}

\textbf{Endpoint:} \texttt{POST /v1/services/transcriber/transcribe}

\textbf{Autenticación:} Requerida

\textbf{Descripción:} Crea o actualiza la transcripción de una lección específica de un curso.

\textbf{Códigos de Respuesta:}
\begin{itemize}
    \item \texttt{201}: Successful Response
    \item \texttt{422}: Validation Error
\end{itemize}

\textbf{Ejemplo:}

\begin{lstlisting}[language=bash]
curl -X POST http://localhost:8000/v1/services/transcriber/transcribe \
  -H "Authorization: Bearer YOUR_ACCESS_TOKEN" \
  -H "Content-Type: application/json" \
  -d '{
    "course_code": "PROG1",
    "lesson_number": 1,
    "model_size": "medium"
  }'
\end{lstlisting}

\subsubsection{Listar transcripciones}
\label{subsubsection:A-3-Listar transcripciones}

\textbf{Endpoint:} \texttt{GET /transcriptions}

\textbf{Autenticación:} Requerida

\textbf{Parámetros de consulta:}
\begin{itemize}
    \item \texttt{course\_code} (opcional): Filtrar por código de curso
\end{itemize}

\textbf{Descripción:} Lista todas las transcripciones disponibles con filtrado opcional.

\textbf{Códigos de Respuesta:}
\begin{itemize}
    \item \texttt{200}: Successful Response
    \item \texttt{422}: Validation Error
\end{itemize}

\textbf{Ejemplo:}

\begin{lstlisting}[language=bash]
# Listar todas las transcripciones
curl -X GET http://localhost:8000/transcriptions \
  -H "Authorization: Bearer YOUR_ACCESS_TOKEN"

# Filtrar por curso
curl -X GET "http://localhost:8000/transcriptions?course_code=PROG1" \
  -H "Authorization: Bearer YOUR_ACCESS_TOKEN"
\end{lstlisting}

\subsection{Gestión de Conversaciones}
\label{subsection:A-3-Gestión de Conversaciones}

\subsubsection{Crear una conversación}
\label{subsubsection:A-3-Crear una conversación}

\textbf{Endpoint:} \texttt{POST /conversations}

\textbf{Autenticación:} Requerida

\textbf{Descripción:} Crea una nueva conversación para el usuario autenticado.

\textbf{Códigos de Respuesta:}
\begin{itemize}
    \item \texttt{201}: Successful Response
    \item \texttt{422}: Validation Error
\end{itemize}

\textbf{Ejemplo:}

\begin{lstlisting}[language=json]
curl -X POST http://localhost:8000/conversations \
  -H "Authorization: Bearer YOUR_ACCESS_TOKEN" \
  -H "Content-Type: application/json" \
  -d '{
    "title": "Dudas sobre Programaci{\'o}n 1"
  }'
\end{lstlisting}

\subsubsection{Listar conversaciones}
\label{subsubsection:A-3-Listar conversaciones}

\textbf{Endpoint:} \texttt{GET /conversations}

\textbf{Autenticación:} Requerida

\textbf{Parámetros de consulta:}
\begin{itemize}
    \item \texttt{limit} (integer, opcional): Límite de resultados por página
    \item \texttt{offset} (integer, opcional): Desplazamiento para paginación
\end{itemize}

\textbf{Descripción:} Recupera una lista paginada de conversaciones del usuario.

\textbf{Códigos de Respuesta:}
\begin{itemize}
    \item \texttt{200}: Successful Response
    \item \texttt{422}: Validation Error
\end{itemize}

\textbf{Ejemplo:}

\# Sin paginación
\begin{lstlisting}[language=bash]
curl -X GET http://localhost:8000/conversations \
  -H "Authorization: Bearer YOUR_ACCESS_TOKEN"
\end{lstlisting}
\# Con paginación
\begin{lstlisting}[language=bash]
curl -X GET "http://localhost:8000/conversations?limit=10&offset=0" \
  -H "Authorization: Bearer YOUR_ACCESS_TOKEN"
\end{lstlisting}

\subsubsection{Obtener mensajes de una conversación}
\label{subsubsection:A-3-Obtener mensajes de una conversación}

\textbf{Endpoint:} \texttt{GET /conversations/\{conversation\_id\}/messages}

\textbf{Autenticación:} Requerida

\textbf{Parámetros de ruta:}
\begin{itemize}
    \item \texttt{conversation\_id} (integer, requerido): ID de la conversación
\end{itemize}

\textbf{Descripción:} Recupera todos los mensajes de una conversación específica del usuario.

\textbf{Códigos de Respuesta:}
\begin{itemize}
    \item \texttt{200}: Successful Response
    \item \texttt{422}: Validation Error
\end{itemize}

\textbf{Ejemplo:}

\begin{lstlisting}[language=bash]
curl -X GET http://localhost:8000/conversations/1/messages \
  -H "Authorization: Bearer YOUR_ACCESS_TOKEN"
\end{lstlisting}

\subsubsection{Renombrar una conversación}
\label{subsubsection:A-3-Renombrar una conversación}

\textbf{Endpoint:} \texttt{PATCH /conversations/\{conversation\_id\}}

\textbf{Autenticación:} Requerida

\textbf{Parámetros de ruta:}
\begin{itemize}
    \item \texttt{conversation\_id} (integer, requerido): ID de la conversación
\end{itemize}

\textbf{Descripción:} Cambia el nombre de una conversación.

\textbf{Códigos de Respuesta:}
\begin{itemize}
    \item \texttt{200}: Successful Response
    \item \texttt{422}: Validation Error
\end{itemize}

\textbf{Ejemplo:}

\begin{lstlisting}[language=bash]
curl -X PATCH http://localhost:8000/conversations/1 \
  -H "Authorization: Bearer YOUR_ACCESS_TOKEN" \
  -H "Content-Type: application/json" \
  -d '{
    "title": "Nuevo nombre de la conversación"
  }'
\end{lstlisting}

\subsubsection{Eliminar una conversación}
\label{subsubsection:A-3-Eliminar una conversación}

\textbf{Endpoint:} \texttt{DELETE /conversations/\{conversation\_id\}}

\textbf{Autenticación:} Requerida

\textbf{Parámetros de ruta:}
\begin{itemize}
    \item \texttt{conversation\_id} (integer, requerido): ID de la conversación
\end{itemize}

\textbf{Descripción:} Elimina una conversación propiedad del usuario autenticado.

\textbf{Códigos de Respuesta:}
\begin{itemize}
    \item \texttt{204}: Successful Response
    \item \texttt{422}: Validation Error
\end{itemize}

\textbf{Ejemplo:}

\begin{lstlisting}[language=bash]
curl -X DELETE http://localhost:8000/conversations/1 \
  -H "Authorization: Bearer YOUR_ACCESS_TOKEN"
\end{lstlisting}

\subsection{Gestión de Cursos}
\label{subsection:A-3-Gestión de Cursos}

\subsubsection{Obtener información de un curso}
\label{subsubsection:A-3-Obtener información de un curso}

\textbf{Endpoint:} \texttt{GET /courses/\{course\_code\}}

\textbf{Autenticación:} Requerida

\textbf{Parámetros de ruta:}
\begin{itemize}
    \item \texttt{course\_code} (string, requerido): Código del curso
\end{itemize}

\textbf{Descripción:} Recupera la información del curso especificado por su código.

\textbf{Códigos de Respuesta:}
\begin{itemize}
    \item \texttt{200}: Successful Response
\end{itemize}

\textbf{Ejemplo:}

\begin{lstlisting}[language=bash]
curl -X GET http://localhost:8000/courses/PROG1 \
  -H "Authorization: Bearer YOUR_ACCESS_TOKEN"
\end{lstlisting}

\subsubsection{Obtener lecciones de un curso}
\label{subsubsection:A-3-Obtener lecciones de un curso}

\textbf{Endpoint:} \texttt{GET /courses/\{course\_code\}/lessons}

\textbf{Autenticación:} Requerida

\textbf{Parámetros de ruta:}
\begin{itemize}
    \item \texttt{course\_code} (string, requerido): Código del curso
\end{itemize}

\textbf{Descripción:} Recupera todas las lecciones de un curso específico.

\textbf{Códigos de Respuesta:}
\begin{itemize}
    \item \texttt{200}: Successful Response
    \item \texttt{422}: Validation Error
\end{itemize}

\textbf{Ejemplo:}

\begin{lstlisting}[language=bash]
curl -X GET http://localhost:8000/courses/PROG1/lessons \
  -H "Authorization: Bearer YOUR_ACCESS_TOKEN"
\end{lstlisting}

\subsubsection{Obtener transcripción de una lección}
\label{subsubsection:A-3-Obtener transcripción de una lección}

\textbf{Endpoint:} \texttt{GET /courses/\{course\_code\}/lessons/\{lesson\_number\}}

\textbf{Autenticación:} Requerida

\textbf{Parámetros de ruta:}
\begin{itemize}
    \item \texttt{course\_code} (string, requerido): Código del curso
    \item \texttt{lesson\_number} (integer, requerido): Número de lección
\end{itemize}

\textbf{Parámetros de consulta:}
\begin{itemize}
    \item \texttt{model\_size} (opcional): Tamaño del modelo de transcripción
\end{itemize}

\textbf{Descripción:} Recupera la transcripción de una lección específica de un curso.

\textbf{Códigos de Respuesta:}
\begin{itemize}
    \item \texttt{200}: Successful Response
    \item \texttt{422}: Validation Error
\end{itemize}

\textbf{Ejemplo:}

\begin{lstlisting}[language=bash]
curl -X GET "http://localhost:8000/courses/PROG1/lessons/1?model_size=medium" \
  -H "Authorization: Bearer YOUR_ACCESS_TOKEN"
\end{lstlisting}

\subsection{Gestión de Usuario}
\label{subsection:A-3-Gestión de Usuario}

\subsubsection{Obtener perfil de usuario}
\label{subsubsection:A-3-Obtener perfil de usuario}

\textbf{Endpoint:} \texttt{GET /users/me}

\textbf{Autenticación:} Requerida

\textbf{Descripción:} Recupera la información del perfil del usuario autenticado actualmente.

\textbf{Códigos de Respuesta:}
\begin{itemize}
    \item \texttt{200}: Successful Response
\end{itemize}

\textbf{Ejemplo:}

\begin{lstlisting}[language=bash]
curl -X GET http://localhost:8000/users/me \
  -H "Authorization: Bearer YOUR_ACCESS_TOKEN"
\end{lstlisting}

